%% Drafted from the mapped corpus by scripts/draft_topics.py.
% Model: gpt-5.6-terra; source characters: 23818.
\begin{konspekt}
\kreq{Дефиниции}{краен ориентиран и краен неориентиран (мулти)граф --- \kref{3.2}.}
\kreq{Дефиниции}{път и цикъл в ориентиран и неориентиран мултиграф; свързаност и свързани компоненти --- \kref{3.3}.}
\kreq{Дефиниции}{дърво, кореново дърво, покриващо дърво --- \kref{3.4}.}
\kreq{Дефиниции}{Ойлеров път и Ойлеров цикъл --- \kref{3.6}.}
\kreq{Доказателства}{всяко кореново дърво е дърво и $|V|=|E|+1$ --- \kref{3.4}.}
\kreq{Доказателство}{теорема за съществуване на Ойлеров цикъл --- \kref{3.6}.}
\kreq{Формулировка}{критерий за съществуване на Ойлеров път --- \kref{3.6}.}
\kreq{Алгоритми}{обхождане в широчина и в дълбочина --- \kref{3.5}.}
\kreq{Задача 1}{покриващо дърво в широчина или в дълбочина --- \kref{3.7.1}.}
\kreq{Задача 2}{Ойлеров цикъл или път, или доказателство, че няма --- \kref{3.7.2}.}
\kreq{Задача 3}{разбиване на ребрата на минимален брой пътища без общо ребро --- \kref{3.7.3}.}
\kprereq{\kr{3.1} повтаря релациите на еквивалентност от Въпрос 1; тук от тях се използва само връзката „достижимост $=$ еквивалентност, класовете $=$ компоненти“.}
\kreq{Литература}{[12], [33], [40].}
\end{konspekt}

\section{Релации на еквивалентност и класове на еквивалентност}\label{s:3.1}

\textbf{Предварителен раздел, който конспектът не изисква по този въпрос.}
Анотацията към Въпрос 3 започва от дефинициите за граф; релациите на
еквивалентност принадлежат на Въпрос 1. Разделът е тук само защото
достижимостта в неориентиран граф е релация на еквивалентност, а класовете й са
свързаните компоненти --- това е използвано по-нататък. При повторение той може
да се прескочи и да се започне направо от \S3.2.

\begin{defn}
Нека $A$ е множество. Релация $\sim\subseteq A\times A$ се нарича \emph{релация на еквивалентност}, ако за всички $x,y,z\in A$ са изпълнени:
\begin{enumerate}
    \item \emph{рефлексивност}: $x\sim x$;
    \item \emph{симетричност}: ако $x\sim y$, то $y\sim x$;
    \item \emph{транзитивност}: ако $x\sim y$ и $y\sim z$, то $x\sim z$.
\end{enumerate}
\end{defn}

\begin{defn}
Нека $\sim$ е релация на еквивалентност върху $A$ и $a\in A$. \emph{Класът на еквивалентност} на $a$ е множеството
\[
[a]_{\sim}=\{x\in A\mid x\sim a\}.
\]
Когато релацията е ясна от контекста, пишем просто $[a]$.
\end{defn}

За да се определят класовете на еквивалентност, за всеки още необработен елемент $a\in A$ се намират всички елементи $x$, за които $x\sim a$. Полученото множество е $[a]$. След това всички негови елементи се отбелязват като обработени и процедурата се повтаря с елемент извън вече намерените класове.

\begin{prop}
Ако $\sim$ е релация на еквивалентност върху $A$, то за произволни $a,b\in A$ класовете $[a]$ и $[b]$ или съвпадат, или са непресичащи се.
\end{prop}

\begin{proof}
Да допуснем, че съществува $x\in [a]\cap[b]$. Тогава $x\sim a$ и $x\sim b$. От симетричността следва $a\sim x$, а от транзитивността с $x\sim b$ получаваме $a\sim b$.

Нека $y\in[a]$. Тогава $y\sim a$ и $a\sim b$, следователно $y\sim b$, т.е. $y\in[b]$. Значи $[a]\subseteq[b]$. Аналогично $[b]\subseteq[a]$, следователно $[a]=[b]$.
\end{proof}

\begin{cor}
Класовете на еквивалентност на релация на еквивалентност образуват разбиение на множеството $A$: всеки елемент принадлежи на точно един клас.
\end{cor}

\begin{proof}
Рефлексивността дава $a\in[a]$, така че класовете покриват $A$. Ако $c\in[a]\cap[b]$, симетричността и транзитивността дават $a\sim b$, а после $x\sim a\Longleftrightarrow x\sim b$. Значи $[a]=[b]$. Следователно различните класове са непресичащи се и всеки елемент е в точно един от тях.
\end{proof}

\begin{example}
Върху $\mathbb{Z}$ да е зададена релацията
\[
a\sim b\Longleftrightarrow 3\mid(a-b).
\]
Тя разделя целите числа на три класа:
\[
[0]=\{\ldots,-6,-3,0,3,6,\ldots\},
\]
\[
[1]=\{\ldots,-5,-2,1,4,7,\ldots\},
\]
\[
[2]=\{\ldots,-4,-1,2,5,8,\ldots\}.
\]
За да се намери класът на дадено число, е достатъчно да се определи остатъкът му при деление на $3$.
\end{example}

\section{Графи и основни понятия}\label{s:3.2}

\subsection{Неориентирани графи}\label{s:3.2.1}

\begin{defn}
Краен неориентиран граф е наредена двойка
\[
G=(V,E),
\]
където $V$ е непразно крайно множество, чиито елементи се наричат \emph{върхове}, а
\[
E\subseteq\{X\subseteq V\mid |X|=2\}
\]
е множество от \emph{ребра}.
\end{defn}

Тъй като всяко ребро е двуелементно множество, в неориентиран граф реброто между $u$ и $v$ се записва като $\{u,v\}$; редът на върховете няма значение. В тази дефиниция няма примки и две еднакви ребра не могат да се срещат.

\begin{defn}
Нека $G=(V,E)$ е граф.
\begin{enumerate}
    \item Ако $E=\varnothing$, графът се нарича \emph{празен}.
    \item Ако $|V|=1$, графът се нарича \emph{тривиален}.
    \item Върховете $u,v\in V$ са \emph{съседни}, ако $\{u,v\}\in E$.
    \item Ребрата $e_1,e_2\in E$ са \emph{инцидентни}, ако
    \[
    e_1\cap e_2\neq\varnothing.
    \]
    \item Реброто $e\in E$ е инцидентно с върха $u\in V$, ако $u\in e$.
\end{enumerate}
\end{defn}

\begin{defn}
За връх $u\in V$ дефинираме:
\[
N(u)=\{v\in V\mid u\text{ и }v\text{ са съседни}\},
\]
\[
N[u]=N(u)\cup\{u\},
\]
\[
I(u)=\{e\in E\mid u\in e\}.
\]
Числото
\[
d(u)=|I(u)|
\]
се нарича \emph{степен} на върха $u$.
\end{defn}

Максималната и минималната степен на графа се означават съответно с
\[
\Delta(G)=\max\{d(u)\mid u\in V\},
\qquad
\delta(G)=\min\{d(u)\mid u\in V\}.
\]

\subsection{Ориентирани графи и мултиграфи}\label{s:3.2.2}

\begin{defn}
Краен ориентиран граф е наредена двойка
\[
G=(V,E),
\]
където $V$ е непразно крайно множество от върхове и
\[
E\subseteq (V\times V)\setminus\{(u,u)\mid u\in V\}.
\]
Елементите на $E$ се наричат ориентирани ребра или дъги.
\end{defn}

Дъгата $(u,v)$ има начало $u$ и край $v$. За разлика от неориентираното ребро, като цяло $(u,v)\neq(v,u)$.

\begin{defn}
Неориентиран мултиграф е наредена тройка
\[
G=(V,E,f_G),
\]
където $V$ е непразно множество от върхове, $E$ е множество от ребра, $V\cap E=\varnothing$, и
\[
f_G:E\longrightarrow\{X\subseteq V\mid |X|=2\}
\]
е свързваща функция.
\end{defn}

Свързващата функция указва краищата на всяко ребро. Различни елементи на $E$ могат да имат един и същ образ чрез $f_G$, поради което мултиграфът позволява успоредни ребра.

\begin{defn}
Ориентиран мултиграф е наредена тройка
\[
G=(V,E,f_G),
\]
където $V$ е непразно множество от върхове, $E$ е множество от ребра, $V\cap E=\varnothing$, и
\[
f_G:E\longrightarrow V\times V
\]
е свързваща функция.
\end{defn}

\begin{supp}[Бележка]
В представените материали за ориентиран мултиграф се допуска $f_G(e)=(u,u)$, тоест примки. Това се различава от дадената по-горе дефиниция за обикновен ориентиран граф, в която примките са изключени.
\end{supp}

\section{Пътища, цикли и свързаност}\label{s:3.3}

\subsection{Пътища и цикли}\label{s:3.3.1}

\begin{defn}
Нека $G=(V,E,f_G)$ е мултиграф. \emph{Път} в $G$ е алтернираща редица от върхове и ребра
\[
p=(u_0,e_{k_0},u_1,e_{k_1},\ldots,e_{k_{l-1}},u_l),
\]
където $l\geq 0$, $u_i\in V$ за $0\leq i\leq l$, $e_{k_j}\in E$ за $0\leq j\leq l-1$, и всяко ребро $e_{k_j}$ свързва последователните върхове $u_j$ и $u_{j+1}$.
\end{defn}

При ориентиран мултиграф изискването е преминаването по всяко ребро да е в посоката, зададена от свързващата функция.

Върховете $u_0$ и $u_l$ се наричат съответно начало и край на пътя, а останалите върхове са вътрешни. Дължината на пътя е броят ребра в него:
\[
\lVert p\rVert=l.
\]

\begin{defn}
Пътят $p$ се нарича \emph{прост}, ако всички негови върхове и ребра са различни.
\end{defn}

\begin{defn}
Пътят
\[
p=(u_0,e_{k_0},u_1,\ldots,e_{k_{l-1}},u_l)
\]
се нарича \emph{цикъл}, ако $u_0=u_l$. Той е \emph{прост цикъл}, ако има поне едно ребро, върховете $u_0,u_1,\ldots,u_{l-1}$ са по двойки различни и ребрата му са по двойки различни. С други думи, единственото допустимо повторение на връх е $u_0=u_l$.
\end{defn}

\subsection{Свързаност}\label{s:3.3.2}

\begin{defn}
Нека $G=(V,E)$ е неориентиран граф. Върховете $u,v\in V$ са \emph{свързани}, ако съществува път от $u$ до $v$. Графът $G$ е \emph{свързан}, ако всеки два негови върха са свързани.
\end{defn}

\begin{defn}
Ориентиран граф е \emph{слабо свързан}, ако неориентираният граф, получен чрез премахване на посоките на дъгите, е свързан. Той е \emph{полусвързан}, ако за всеки два върха $u$ и $v$ има ориентиран път поне в едната посока: от $u$ до $v$ или от $v$ до $u$. Той е \emph{силно свързан}, ако за всеки два върха има ориентиран път и от първия до втория, и от втория до първия.
\end{defn}

\begin{defn}
Нека $G=(V,E)$ е неориентиран граф. Релацията на достижимост $R_G$ върху $V$ е зададена чрез
\[
R_G\subseteq V\times V,
\qquad
uR_Gv\Longleftrightarrow u\text{ и }v\text{ са свързани}.
\]
\end{defn}

\begin{prop}
Релацията $R_G$ на достижимост в неориентиран граф е релация на еквивалентност.
\end{prop}

\begin{proof}
За всеки връх $u$ съществува тривиален път с дължина $0$ от $u$ до $u$, следователно $uR_Gu$.

Ако има път от $u$ до $v$, същата последователност от ребра, обходена в обратен ред, е път от $v$ до $u$. Следователно $uR_Gv$ предполага $vR_Gu$.

Накрая, ако има път от $u$ до $v$ и път от $v$ до $w$, те могат да се съединят в път от $u$ до $w$. Следователно релацията е транзитивна.
\end{proof}

\begin{defn}
Подграфите на $G$, индуцирани от класовете на еквивалентност на $R_G$, се наричат \emph{свързани компоненти} на $G$.
\end{defn}

Еквивалентно, свързаните компоненти са максималните по включване свързани подграфи. Следователно намирането на свързаните компоненти е конкретен случай на определяне на класовете на еквивалентност на релация.

\section{Дървета}\label{s:3.4}

\begin{defn}
\emph{Дърво} е неориентиран граф, който е свързан и ацикличен.
\end{defn}

Тривиалният граф е дърво. В дървото между два върха няма цикъл; при добавяне на нов връх към дърво и свързването му с точно едно ребро към стар връх не се създава цикъл.

\begin{defn}
\emph{Кореново дърво} определяме индуктивно. Базата е тривиалният граф
\[
(\{r\},\varnothing),
\]
който има корен $r$. В индуктивната стъпка към вече построено кореново дърво се добавя нов връх $u$ и едно ребро, което го свързва с връх $v$ от старото дърво. Коренът се запазва.
\end{defn}

\begin{thm}
Всяко кореново дърво е дърво.
\end{thm}

\begin{proof}
Доказателството е чрез индукция по построението. Базовият едновърхов граф е свързан и няма цикъл.

Нека $T=(V,E)$ е свързано и ациклично дърво. Добавяме нов връх $u\notin V$ и едно ребро $\{v,u\}$, където $v\in V$. Полученият граф е свързан, защото от всеки стар връх има път до $v$, а след това реброто $\{v,u\}$ води до $u$. Новият връх е инцидентен само с едно ребро, затова не може да участва в цикъл. Старите ребра не образуват цикъл по индукционното предположение. Следователно полученият граф също е дърво.
\end{proof}

Обратно, всяко крайно дърво с избран корен може да се получи чрез това построение. Ако има повече от един връх, дървото има листо, различно от корена. Премахването на листото и единственото му инцидентно ребро оставя дърво със същия корен; твърдението следва с индукция по броя върхове, като накрая листото и реброто се добавят обратно.

Следователно кореновото дърво може еквивалентно да се разглежда като дърво с избран връх $r$, наречен \emph{корен}. Изборът на корен задава йерархия: всеки връх, различен от корена, има единствен непосредствен предшественик по пътя към корена.

\begin{thm}
Нека $T=(V,E)$ е кореново дърво. Тогава
\[
|V|=|E|+1.
\]
\end{thm}

\begin{proof}
Използваме индукция по построението на кореновото дърво. За базата
\[
T=(\{r\},\varnothing)
\]
имаме $|V|=1=0+1=|E|+1$.

При индуктивната стъпка се добавят едновременно един нов връх и едно ново ребро. Ако за старото дърво е изпълнено $|V|=|E|+1$, то за новото дърво $T'$ е в сила
\[
|V(T')|=|V|+1=|E|+2=|E(T')|+1.
\]
\end{proof}

\begin{defn}
Нека $G=(V,E)$ е граф. \emph{Покриващо дърво} на $G$ е дърво
\[
T=(V,E'),
\qquad E'\subseteq E.
\]
\end{defn}

Покриващото дърво съдържа всички върхове на графа, но само част от неговите ребра.

\begin{thm}
Краен непразен неориентиран граф $G=(V,E)$ има покриващо дърво тогава и само тогава, когато е свързан.
\end{thm}

\begin{proof}
Ако има покриващо дърво, пътят в него между всяка двойка върхове е път и в $G$, следователно $G$ е свързан. Обратно, за краен непразен свързан неориентиран граф започваме с един връх. Докато не сме включили всички върхове, свързаността осигурява ребро от включен връх към невключен. Добавяме това ребро и новия връх. Свързаността на построеното подмножество се запазва и цикъл не възниква, защото новият връх има само едно добавено ребро. След $|V|-1$ стъпки получаваме покриващо дърво.
\end{proof}

\section{Обхождания на графи}\label{s:3.5}

Обхождането систематично посещава върховете, достижими от начален връх. При откриването на нов връх се запомня реброто, по което е достигнат; така получените ребра образуват покриващо дърво на достигнатата свързана компонента.

\subsection{Стек и опашка}\label{s:3.5.1}

\begin{defn}
\emph{Стек} е структура от тип LIFO: последно добавеният елемент се извежда първи. Празният стек означаваме с $\{\}$. Ако $S$ е стек и $x$ е елемент, то $\operatorname{push}(S,x)$ добавя $x$ на върха на стека, $\operatorname{top}(S)$ връща горния елемент, а $\operatorname{pop}(S)$ премахва горния елемент.
\end{defn}

\begin{defn}
\emph{Опашка} е структура от тип FIFO: първо добавеният елемент се извежда първи. Операцията $\operatorname{push}(Q,x)$ добавя елемент в края, $\operatorname{top}(Q)$ връща първия елемент, а $\operatorname{pop}(Q)$ го премахва.
\end{defn}

\subsection{Схема за обхождане}\label{s:3.5.2}

Нека върховете на ориентиран мултиграф са означени с $1,\ldots,n$, а началният връх е $1$. Поддържа се булев масив $\operatorname{visited}[1,n]$. В началото всички стойности са \texttt{False}; началният връх се поставя в помощна структура $S$ и се маркира като посетен.

Докато $S$ не е празна, от нея се извежда връх $x$. За всяка дъга $(x,y)$, ако $y$ не е посещаван, той се маркира като посетен и се добавя в $S$. Изборът на конкретна структура за $S$ определя вида обхождане.

\subsection{Обхождане в широчина}\label{s:3.5.3}

\begin{keydefn}
\emph{Обхождане в широчина} (BFS) е обхождане по общата схема, при което помощната структура е опашка.
\end{keydefn}

BFS първо разглежда всички върхове на разстояние $1$ от началния връх, след това върховете на разстояние $2$ и т.н. Ако при откриването на $y$ чрез $x$ запомним $\pi[y]=x$, стойностите $\pi$ задават дърво на обхождането с корен началния връх.

\begin{verbatim}
BFS(G=(V,E))
for i <- 1 to n
    colour[i] <- white
    pi[i] <- NIL
    d[i] <- infinity
colour[1] <- gray
d[1] <- 0
create queue Q
Enqueue(Q,1)
while not isempty(Q)
    x <- Dequeue(Q)
    for y in adj(x)
        if colour[y] = white
            colour[y] <- gray
            pi[y] <- x
            d[y] <- d[x] + 1
            Enqueue(Q,y)
    colour[x] <- black
return colour,pi,d
\end{verbatim}

\begin{keythm}
Обхождането в широчина намира най-късите пътища от началния връх до всички достижими върхове в графа.
\end{keythm}

\begin{proof}
Разстоянието тук е броят на ребрата. Началният връх $s$ получава $d[s]=0$. FIFO опашката обработва изцяло върховете на слой $k$, преди върховете на слой $k+1$. По индукция нека всички върхове с действително разстояние най-много $k$ са открити с правилни стойности. Връх $v$ на разстояние $k+1$ има предшественик на разстояние $k$ по най-къс път; при обработката му $v$ се открива най-късно със стойност $k+1$. Всяка присвоена стойност е дължина на реален път и не може да е по-малка от най-късото разстояние. Така $d[v]=k+1$. Всички достижими върхове попадат в някой слой, което завършва доказателството.
\end{proof}

Стойността $d[y]$, присвоена при първото откриване на $y$, е дължината на най-къс път от началния връх до $y$. Това следва от факта, че опашката обработва върховете по нарастване на разстоянието от началния връх.

\subsection{Обхождане в дълбочина}\label{s:3.5.4}

\begin{keydefn}
\emph{Обхождане в дълбочина} (DFS) е обхождане, при което се продължава към необходен съсед, докато това е възможно; когато няма такъв съсед, алгоритъмът се връща назад. Итеративната реализация използва стек, тоест структура LIFO.
\end{keydefn}

Следната рекурсивна реализация изразява същата идея:

\begin{verbatim}
DFS(G=(V,E))
for i <- 1 to n
    colour[i] <- white
    N[i] <- NIL
DFS-Visit(G,1)
return colour,N

DFS-Visit(G=(V,E), x in V)
colour[x] <- gray
foreach y in adj[x]
    if colour[y] = white
        N[y] <- x
        DFS-Visit(G,y)
colour[x] <- black
\end{verbatim}

При достигане на необходен връх $y$ от $x$, стойността $N[y]=x$ запомня непосредствения предшественик на $y$ в дървото на DFS. Ако графът не е свързан, едно обхождане от избран начален връх посещава само неговата компонента; за обхождане на всички върхове алгоритъмът се стартира от всеки все още необходен връх.

И DFS, и BFS могат да бъдат изменени за намиране на път между два върха: единият се избира за начален, а обхождането се прекратява при достигане на другия. Пътят се възстановява чрез запомнените предшественици.

\section{Ойлерови пътища и цикли}\label{s:3.6}

\begin{defn}
Нека $G$ е свързан мултиграф. \emph{Ойлеров път} е път, който преминава точно веднъж през всяко ребро на графа. \emph{Ойлеров цикъл} е Ойлеров път, чийто начален и краен връх съвпадат. Граф, който притежава Ойлеров цикъл, се нарича \emph{Ойлеров граф}.
\end{defn}

\begin{keythm}[Критерий за Ойлеров цикъл в свързан мултиграф]
Нека $G$ е свързан мултиграф (в частност --- свързан граф) с поне едно ребро. Тогава $G$ е Ойлеров тогава и само тогава, когато всеки негов връх има четна степен. Кратните ребра и примките не изискват отделно разглеждане: всяка примка добавя $2$ към степента на върха си, а кратните ребра се броят поотделно.
\end{keythm}

\begin{proof}
Нека $G$ има Ойлеров цикъл. Всеки път, когато цикълът влиза във връх, той трябва да го напусне по друго ребро. Следователно ребрата, инцидентни с произволен връх, се групират по двойки: едно за влизане и едно за излизане. Значи степента на всеки връх е четна.

Обратно, нека $G$ е свързан и всеки връх има четна степен. Избираме начален връх и вървим по неизползвани ребра, като всяко използвано ребро се маркира. Пътят не може да завърши в връх, различен от началния: при такъв краен връх броят използвани инцидентни ребра би бил нечетен, което противоречи на четността на степента му. Така получаваме цикъл.

Ако цикълът съдържа всички ребра, той е Ойлеров. В противен случай поради свързаността съществува връх от получения цикъл, инцидентен с неизползвано ребро. От този връх се построява нов цикъл по неизползвани ребра и се вмъква в стария. Процесът се повтаря. При всяка стъпка броят неизползвани ребра намалява, затова след краен брой стъпки се получава цикъл, съдържащ всяко ребро точно веднъж.
\end{proof}

\begin{cor}
Свързан мултиграф има Ойлеров път, който не е цикъл, тогава и само тогава, когато точно два негови върха са с нечетна степен, а всички останали са с четна степен.
\end{cor}

\begin{proof}
При Ойлеров път всяко посещение на вътрешен връх използва двойка ребра: за влизане и излизане. Само началото и краят имат по едно несдвоено ребро, затова са точно двата върха с нечетна степен. Обратно, ако нечетните върхове са $u,v$, добавяме още едно ребро между тях. Всички степени стават четни, а графът остава свързан, затова по теоремата за Ойлеров цикъл има такъв цикъл. Премахваме добавеното ребро и прочитаме останалата част от цикъла от единия му край: получаваме Ойлеров път от $u$ до $v$.
\end{proof}

\section{Примерни задачи по анотацията}\label{s:3.7}

Анотацията назовава три типа задачи. Общите алгоритми са дадени по-горе; тук и
трите са изпълнени докрай върху конкретни входове.

Навсякъде в този раздел работим с графа
\[
G:\quad V=\{1,2,3,4,5,6\},\qquad
E=\{12,\,13,\,23,\,24,\,35,\,45,\,46,\,56\},
\]
където $uv$ означава реброто между $u$ и $v$. Списъците на съседите са подредени
по нарастване, за да е обхождането еднозначно:
\[
1:2,3\quad 2:1,3,4\quad 3:1,2,5\quad 4:2,5,6\quad 5:3,4,6\quad 6:4,5 .
\]
Степените са
\[
\deg 1=2,\quad \deg 2=3,\quad \deg 3=3,\quad
\deg 4=3,\quad \deg 5=3,\quad \deg 6=2 .
\]

\subsection{Задача 1: покриващо дърво в широчина и в дълбочина}\label{s:3.7.1}

\begin{example}[BFS от връх $1$]
Опашката се обработва отляво надясно; при всеки връх се обхождат необходените му
съседи по реда от списъка.

\begin{center}
\begin{tabular}{llll}
\toprule
Изваден връх & Нови върхове & Ребра на дървото & Опашка след стъпката\\
\midrule
$1$ & $2,3$ & $12,\,13$ & $2,3$\\
$2$ & $4$   & $24$      & $3,4$\\
$3$ & $5$   & $35$      & $4,5$\\
$4$ & $6$   & $46$      & $5,6$\\
$5$ & ---   & ---       & $6$\\
$6$ & ---   & ---       & ---\\
\bottomrule
\end{tabular}
\end{center}

Покриващото дърво е $\{12,13,24,35,46\}$ --- пет ребра при шест върха, както
изисква $|V|=|E|+1$. Разстоянията от началния връх са
\[
d(1)=0,\quad d(2)=d(3)=1,\quad d(4)=d(5)=2,\quad d(6)=3,
\]
и това са дължините на най-късите пътища, защото BFS обхожда върховете по
нарастващо разстояние.
\end{example}

\begin{example}[DFS от връх $1$]
Рекурсията влиза колкото може по-дълбоко, преди да се върне:
\[
1\to 2\to 3\to 5\to 4\to 6,
\]
след което от $6$ няма необходени съседи и обхождането се връща по стека до $1$.
Записаните предшественици са
\[
N[2]=1,\quad N[3]=2,\quad N[5]=3,\quad N[4]=5,\quad N[6]=4,
\]
тоест покриващото дърво е $\{12,23,35,45,46\}$ и е път. Ребрата $13$, $24$ и
$56$ остават извън дървото --- те са обратни ребра към вече посетени върхове.
\end{example}

\begin{remark}
Двете дървета са различни, макар графът и началният връх да са едни и същи. При
такава задача винаги се посочва и редът на съседите: без него отговорът не е
еднозначен.
\end{remark}

\subsection{Задача 2: Ойлеров цикъл или доказателство, че такъв няма}\label{s:3.7.2}

\begin{example}[Отрицателен отговор за $G$]
В $G$ върховете $2,3,4,5$ са с нечетна степен. По критерия Ойлеров цикъл няма,
защото не всички степени са четни; по следствието няма и Ойлеров път, защото
върховете с нечетна степен не са точно два, а четири. Това е пълен отговор ---
позоваване на критерия, не търсене на цикъл „на ръка“.
\end{example}

\begin{example}[Построяване по Hierholzer]
Нека $K_5$ е пълният граф върху $\{1,2,3,4,5\}$. Всички степени са $4$, графът е
свързан, следователно има Ойлеров цикъл. Строим го на три стъпки.

Първо тръгваме от $1$ по неизползвани ребра, докато се върнем в $1$:
\[
C_0:\ 1\to 2\to 3\to 1
\qquad(\text{ребра } 12,23,13).
\]
Връх $2$ още има неизползвани ребра. Строим цикъл от $2$:
\[
2\to 4\to 5\to 2
\qquad(\text{ребра } 24,45,25),
\]
и го вмъкваме в $C_0$ на мястото на $2$:
\[
C_1:\ 1\to 2\to 4\to 5\to 2\to 3\to 1 .
\]
Остават $14,15,34,35$. Строим цикъл от $3$:
\[
3\to 4\to 1\to 5\to 3
\qquad(\text{ребра } 34,14,15,35),
\]
и го вмъкваме на мястото на $3$:
\[
C_2:\ 1\to 2\to 4\to 5\to 2\to 3\to 4\to 1\to 5\to 3\to 1 .
\]
Използвани са всички $\binom{5}{2}=10$ ребра, всяко по веднъж, а началото
съвпада с края. Това е Ойлеров цикъл.
\end{example}

\subsection{Задача 3: разбиване на ребрата на минимален брой пътища без общо ребро}\label{s:3.7.3}

Тук „път“ се разбира в смисъла на \S3.3: върхове могат да се повтарят. Иска се
разбиване на $E$, тоест всяко ребро принадлежи на точно един от пътищата и се
използва в него точно веднъж.

\begin{keythm}[Минимален брой пътища]
Нека $G$ е свързан мултиграф с поне едно ребро и нека точно $2k$ от върховете му
са с нечетна степен (броят на нечетните върхове винаги е четен). Тогава
минималният брой пътища, на които може да се разбие $E$ така, че никои два да
нямат общо ребро, е
\[
\max(k,1).
\]
\end{keythm}

\begin{proof}
\emph{Долна граница.} Разглеждаме произволно разбиване на $E$ на пътища и
фиксираме връх $v$. Всяко преминаване на път през $v$ като вътрешен връх
използва две от инцидентните на $v$ ребра; същото важи и когато път започва и
завършва в $v$. Само път, за който $v$ е точно единият от двата края, използва
нечетен брой инцидентни на $v$ ребра. Тъй като ребрата са разпределени между
пътищата без повторение,
\[
\deg v\equiv\#\{\text{пътища, за които } v \text{ е точно единият край}\}\pmod 2 .
\]
Следователно всеки връх с нечетна степен е край на поне един път. Всеки път има
най-много два края, а нечетните върхове са $2k$, откъдето пътищата са поне $k$.
Освен това $E\ne\varnothing$, тоест поне един път е нужен винаги.

\emph{Конструкция, достигаща границата.} Ако $k=0$, всички степени са четни и по
критерия за Ойлеров цикъл целият граф е един-единствен път (при това затворен),
тоест $1$ път е достатъчен.

Нека $k\geq 1$ и нека нечетните върхове са $u_1,v_1,\ldots,u_k,v_k$ --- $2k$
различни върха, разбити произволно на $k$ двойки. Добавяме $k$ нови ребра
$u_iv_i$. Всяка степен нараства с точно $1$ при нечетните върхове и остава
непроменена при останалите, затова в получения мултиграф $G'$ всички степени са
четни; $G'$ остава свързан. По критерия $G'$ има Ойлеров цикъл $C$.

Премахваме от $C$ добавените $k$ ребра. Понеже $2k$-те върха са различни, никои
две добавени ребра нямат общ връх, следователно никои две не стоят една до друга
в $C$; значи между всеки две последователни премахнати ребра остава поне едно
първоначално ребро. Циклично $C$ се разпада на точно $k$ непразни дъги. Всяка
дъга е път, дъгите нямат общо ребро, а обединението им е точно $E$. Получихме
$k$ пътя.
\end{proof}

\begin{cor}
За граф с няколко свързани компоненти минималният брой пътища е сумата от
$\max(k_i,1)$ по компонентите $i$, съдържащи поне едно ребро. Изолираните
върхове не изискват път.
\end{cor}

\begin{proof}
Път е свързан подграф, затова никой път не пресича две компоненти. Задачата се
разпада независимо по компонентите и се прилага теоремата за всяка от тях.
\end{proof}

\begin{example}[Разбиване на $G$]
В $G$ нечетните върхове са $2,3,4,5$, тоест $2k=4$ и $k=2$: два пътя са
необходими и достатъчни.

Двойките избираме кръстосано --- $(2,5)$ и $(3,4)$ --- защото тези ребра още не
съществуват. Добавяме $e_a=25$ и $e_b=34$. Сега всички степени са четни:
\[
\deg 1=2,\ \deg 2=4,\ \deg 3=4,\ \deg 4=4,\ \deg 5=4,\ \deg 6=2 .
\]
Един Ойлеров цикъл в разширения мултиграф е
\[
2\xrightarrow{12}1\xrightarrow{13}3\xrightarrow{e_b}4\xrightarrow{24}2
\xrightarrow{23}3\xrightarrow{35}5\xrightarrow{56}6\xrightarrow{46}4
\xrightarrow{45}5\xrightarrow{e_a}2 .
\]
Премахваме $e_a$ и $e_b$ и оставащите две дъги са търсените пътища:
\[
P_1:\ 2-1-3,
\qquad
P_2:\ 4-2-3-5-6-4-5 .
\]
Проверка: $P_1$ използва $12,13$; $P_2$ използва $24,23,35,56,46,45$. Заедно това
са всичките $8$ ребра, всяко по веднъж, и общо ребро няма. Краищата на двата
пътя са $2,3$ и $4,5$ --- точно четирите нечетни върха, както изисква долната
граница.
\end{example}

\begin{remark}
Отговорът не е единствен: друг избор на двойки или друг Ойлеров цикъл дава друга
двойка пътища. Единствен е \emph{броят} им. На изпита се иска и обосновка на
минималността, тоест аргументът за долната граница чрез нечетните върхове.
\end{remark}

\section{Изпитен фокус}\label{s:3.8}

\begin{itemize}
    \item Да се дадат определения за неориентиран граф, ориентиран граф, мултиграф, съседство, инцидентност, степен на връх, път, прост път и цикъл.
    \item Да се дефинират свързан граф, слаба и силна свързаност, релация на достижимост и свързани компоненти.
    \item Да се възпроизведе доказателствената идея, че достижимостта в неориентиран граф е релация на еквивалентност; класовете $\,$ѝ са свързаните компоненти.
    \item Да се дефинират дърво, кореново дърво и покриващо дърво.
    \item Да се докаже чрез индукция, че за кореново дърво е изпълнено $|V|=|E|+1$.
    \item Да се формулира критерият: граф има покриващо дърво точно когато е свързан.
    \item Да се опишат BFS и DFS, ролята на опашката при BFS и на стека или рекурсията при DFS, както и дървото на обхождането.
    \item Да се формулира, че BFS намира най-къси пътища от началния връх.
    \item Да се дефинират Ойлеров път и Ойлеров цикъл.
    \item Да се формулира и обясни критерият за Ойлеров цикъл: свързан граф е Ойлеров точно когато всички върхове са с четна степен.
    \item Да се опише идеята на алгоритъма на Hierholzer: строене на цикъл по неизползвани ребра и вмъкване на допълнителни цикли.
    \item Да се формулира условието за Ойлеров път, който не е цикъл: точно два върха са с нечетна степен.
    \item \textbf{Задача 1 от анотацията:} да се построи покриващо дърво на зададен граф в широчина или в дълбочина, с изписан ред на обхождането и на ребрата на дървото (\S3.7.1).
    \item \textbf{Задача 2 от анотацията:} да се построи Ойлеров цикъл или път в зададен мултиграф чрез вмъкване на цикли, или да се докаже, че такъв няма, чрез броя на върховете с нечетна степен (\S3.7.2).
    \item \textbf{Задача 3 от анотацията:} да се разбият ребрата на минимален брой пътища без общо ребро --- долна граница $k$ от $2k$-те нечетни върха, конструкция чрез сдвояване, Ойлеров цикъл и премахване на добавените ребра (\S3.7.3).
    \item Релациите на еквивалентност в \S3.1 са предварителен материал от Въпрос 1; по този въпрос от тях се използва само, че достижимостта е еквивалентност и класовете й са свързаните компоненти.
\end{itemize}
